% Compile twice: lualatex 20001197.tex
\documentclass[11pt]{article}
\usepackage[margin=25mm]{geometry}
\usepackage{fontspec,unicode-math}
\setmainfont{Latin Modern Roman}
\setmathfont{Latin Modern Math}
\usepackage{amsmath,amssymb,mathtools}
\usepackage{xurl,hyperref}
\hypersetup{colorlinks=true,urlcolor=blue,linkcolor=blue}
\setcounter{secnumdepth}{0}
\setlength{\parindent}{0pt}
\setlength{\parskip}{5pt}
\setlength{\emergencystretch}{3em}
\newcommand{\sref}[2]{\hyperlink{source:#1:#2}{[S#2]}}
\newcommand{\cataloguescope}[1]{\paragraph{Source question.}#1}
\newcommand{\floor}[1]{\left\lfloor#1\right\rfloor}
\begin{document}
% UnsolvedMath ID 20001197; source catalogue code AIM-COMPUTATION-0035
\setcounter{section}{20001196}
\section{Exact CI rationalization through rational local charts}\label{20001197}
{\small\sffamily UnsolvedMath ID 20001197 · Computer Science}
\subsection{Definitions and mathematical statement}

\paragraph{Shared notation.}Standard number systems are denoted by $\mathbb N,\mathbb Z,\mathbb Q,\mathbb R,\mathbb C$; index ranges, computational models and specialized notation follow the cited problem source and the definitions below.


A conditional-independence (CI) structure records exactly which relations $X_A\perp X_B\mid X_C$ hold for disjoint index sets. It is probabilistically representable when it is the CI structure of some random vector. A finite uniform probability space gives each elementary outcome the same probability, while the random vector may identify multiple elementary outcomes.
\paragraph{Statement recorded in the supplied source.}
Conjecture 21. Every p-representable CI-structure can be p-represented by a random vector living on a finite probability space endowed with the uniform distribution. \sref{20001197}{2}
\paragraph{Scope qualification.}
The source uses p-representable for probabilistic representability; uniformity concerns elementary outcomes, not necessarily the distinct values of the random vector.
\paragraph{Source status.}
The supplied catalogue labels this record partially solved. Its source-status metadata was checked against the supplied snapshot on 2026-10-11. This import does not independently establish current open status or mathematical correctness.
\subsection{Short English statement}
Can every probabilistically realizable conditional-independence pattern be realized on a finite uniform probability space?
\subsection{Sources}
\begin{itemize}
\item[\hypertarget{source:20001197:1}{S1}] Ulam AI and the UnsolvedMath Contributors, UnsolvedMath: A Curated Collection of Open Mathematics Problems, record 20001197 (AIM-COMPUTATION-0035). Catalogue curation and metadata: CC BY 4.0. The supplied source question is retained; definitions and the English summary are editorial additions.
\url{https://huggingface.co/datasets/ulamai/UnsolvedMath}.
\item[\hypertarget{source:20001197:2}{S2}] AIM Workshop Problem Lists, original problem source linked by catalogue record AIM-COMPUTATION-0035
\url{https://aimath.org/WWN/compalgstat/compalgstat.pdf}.
\end{itemize}
\paragraph{Preparation and review.}
Original question attributed to the cited source and UnsolvedMath contributors. Codex (GPT-6; exact serving revision not exposed) assisted with definitions, English summaries, source inspection and formatting. No new human mathematical authorship or independent proof verification is asserted. Canonical ID 20001197 and source code AIM-COMPUTATION-0035 checked against the complete supplied registry and existing repository statements; known flagged repetitions and obvious identical conjectures excluded. Full mathematical-equivalence/literature review has not been performed. Statement preparation and final review followed contribution prompts C and D; automated checks verify structure and identity, not open status.
\label{end:20001197}
\end{document}
